\section{Extension classes and normalized ramification charts}
\label{sec:local-algebra-v167}
We record two local consequences that separate properties of the
embedded curve from properties of its parameter space.

\subsection{The presentation complex and explicit splitting}
For a polynomial presentation $R=P/J$ and an $R$-module $N$,
square-zero $\C$-algebra extensions with specified kernel $N$
and specified quotient $R$ are classified by
$\operatorname{Ext}^1_R(L_{R/\C},N)$. We use cohomological degrees:
$J/J^2$ is in degree $-1$ and
$\Omega_{P/\C}\otimes_PR$ in degree $0$. The group is
\[
 \coker\bigl(\Hom_R(\Omega_{P/\C}\otimes_PR,N)
                       \longrightarrow\Hom_R(J/J^2,N)\bigr).
\]
Indeed a lift of polynomial generators gives a conormal map;
changing lifts changes it by a derivation. The part of the
cotangent complex in degrees at most $-2$ contributes only in
Ext degrees at least two when the target is in degree zero.
This justifies truncation for this particular group, not a claim
that the whole cotangent complex is two-term.

In the chart of Theorem~\ref{thm:extension-class-v166}, invert $C$.
Then the extension algebra is explicitly
\[
 \C[\lambda,A,B,C,C^{-1},n]/(n^2,An,Bn),\qquad e=n/C.
\]
Writing its elements as $r+an$ gives multiplication
$(r+an)(s+bn)=rs+(rb+sa)n$, where the coefficients of $n$ are
restricted modulo $(A,B)$. This is the actual split algebra on
this chart. On an overlap of two obstruction charts, let $\xi_i$
be the conormal representative obtained from the same algebra
extension. A change of polynomial lifts changes $\xi_i$ by a
derivation. Therefore $[\xi_i]=[\xi_j]$ under the induced overlap
map on the Ext sheaves. Taking the trivialization that sends
$1$ to $[\xi_i]$ makes the overlap function exactly $1$, rather
than an unspecified unit. This proves the asserted trivialization
of the local obstruction sheaf. It does not trivialize the
nilradical line bundle on the singular-conic plane, and it does
not identify the entire global hyper-Ext group with its sheaf of
local degree-one groups. The singular-conic plane is the locus
where the linear part of the conic equation at the attaching
point vanishes; the doubled-line curve is its discriminant-zero
locus. Thus the support of the embedded associated subvariety
and the obstruction to an algebra retraction are different loci.

\begin{proposition}[Singularities and divisor classes after ramification]
\label{prop:toric-classes-v167}
Let $m\geq1$ and put $g=\gcd(m,2)$. For the normalized transverse
surface in Theorem~\ref{thm:ramified-v166}, the character lattice
and its saturated semigroup are
\[
 M_m=\{(a,b)\in\mathbb Z^2:a-2b\equiv0\pmod m\},\qquad
                         M_m\cap\mathbb R_{\geq0}^2.
\]
The primitive rays of the dual cone in $N_m=\Hom(M_m,\mathbb Z)$
are $(1/g,0)$ and $(0,1)$. Its divisor class group is
$\mathbb Z/(m/g)$. For odd $m$ it is the cyclic quotient of type
$\frac1m(1,-2)$; it is smooth for $m=1$ and has a unique singular
point for every odd $m>1$. For even $m=2q$ it is of type
$\frac1q(1,-1)$, with equation $eh=\tau^q$; it is smooth for
$q=1$ and has a unique singular point otherwise. In the full
chart the additional coordinate $\lambda$ gives the product with
an affine line. The singular locus is correspondingly empty or
that line.
\end{proposition}
\begin{proof}
The exponent lattice and its saturation were computed in
Theorem~\ref{thm:ramified-v166}. Their dual lattice is
$N_m=\mathbb Z^2+\mathbb Z(1/m,-2/m)$. On the first coordinate
axis its primitive nonzero vector is $(1/g,0)$; on the second it
is $(0,1)$. This follows by solving the congruence in the unused
coordinate. The index of the lattice generated by these two rays
in $N_m$ is $m/g$. The torus-invariant divisor sequence is the
cokernel of
\[
 M_m\longrightarrow\mathbb Z^2,\qquad (a,b)\longmapsto(a/g,b).
\]
For odd $m$ its image is $a-2b\equiv0\pmod m$. For even
$m=2q$ write $a=2a'$; its image is
$a'-b\equiv0\pmod q$. In both cases the cokernel is the stated
cyclic group. This computes the class group of the normal affine
toric surface, including the smooth cases.

For odd $m$ neither weight has a nontrivial stabilizer on a
coordinate axis away from the origin. The quotient is therefore
smooth away from the origin. Its cone is unimodular exactly when
$m=1$, proving singularity at the origin otherwise. For even
$m$ taking invariants under the order-two pseudoreflection first
gives coordinates $s^2,r$ with residual opposite weights of order
$q$. The displayed hypersurface and its derivatives prove the
remaining assertions. Formation of the product with the
independent coordinate does not change this transverse analysis.
\end{proof}

\begin{lemma}[Gluing integral closures]\label{lem:normalization-gluing-v167}
Let an integral Noetherian scheme be covered by affine charts
$\Spec A_i$ in a fixed function field. Their integral closures
in that field glue to its normalization. In particular the
ramified invariant charts and the unchanged primitive charts
in Theorem~\ref{thm:ramified-v166} glue without a choice of a
cyclic cover or an ordering of branches.
\end{lemma}
\begin{proof}
Integral closure commutes with localization. One implication is
immediate by localizing a monic equation. For the other, clear
the finitely many denominators in a monic equation for an element
integral over a localization: a suitable multiple of that element
by an element of the multiplicative set is integral over the
original ring. Thus both chart normalizations restrict to the
same ring on every principal overlap. Affine refinements and the
cocycle identity of localization give the gluing. All rings here
are finite-type over $\C$, so these normalizations are finite.
\end{proof}
